\section{Introduction}
Sports facility aggregators sell capacity that expires. An empty court or pitch cannot recover the revenue from a slot once its start time has passed. Demand varies with time, weather, facility type and price. The research problem is how to forecast booking probability and use that forecast to choose prices within constraints set by facility owners.

The literature review and original Methods proposed an integrated forecasting and pricing study using real aggregator records. The present article reports a preliminary simulation using the supplied booking generator. The simulated records permit evaluation of forecasts and pricing calculations against a known booking mechanism. They do not provide evidence about the behaviour of actual customers or the revenue of an operating aggregator.

The review asked which forecasting methods and external variables are useful, how pricing research represents price response and capacity, and which gaps remain for bookable sports facilities. These questions lead to three operational questions for the experiment. RQ1 asks whether the selected forecasting model improves probabilistic accuracy over the baselines. RQ2 asks what revenue change the constrained pricing policy produces in the simulation, including its uncertainty and agreement with the simulator. RQ3 asks whether that policy increases weekday off-peak occupancy without an unacceptable increase in cancellations. The success criteria proposed in Methods are retained, subject to the deviations stated below.

\section{Literature review}
The supplied review contains 13 original studies and 10 review papers published from 2021 to 2025. Its search covered scholarly journal sources and publisher or institutional records from ScienceDirect, Springer Nature, INFORMS, Wiley, ACM, MDPI and related indexes. The selection excluded conference proceedings, book chapters and unpublished preprints. The methodological references added in the Methods section serve a different purpose and are not counted as part of this review corpus.


\subsection{Pricing and marketplace mechanisms}

\subsubsection{Surge pricing and two-sided temporal responses}
Hu, Hu, and Zhu study surge pricing in ride-hailing marketplaces and focus on the fact that riders and drivers respond to prices at different speeds \cite{r2}. The problem is therefore not simply choosing a price but understanding the temporal interaction between two sides of a platform. Their methodology uses a two-period game-theoretic model with strategic consumer and driver responses. The study identifies shortlived high-price ``skimming'' surge and low-initial/high-later ``penetration'' surge equilibria, with the latter depending on information sharing with drivers. The contribution is a formal demonstration that the timing of behavioral responses changes the performance of dynamic pricing. A major strength is the explicit two-sided temporal structure; a limitation is that drivers can relocate, whereas sports facilities are fixed assets. For sports facilities, the result suggests that customer reactions to price should be modeled over time rather than assumed to be instantaneous.

\subsubsection{Driver-side surge incentives}
Garg and Nazerzadeh examine how surge pricing affects driver incentives in ride-hailing \cite{r3}. The problem is that a price multiplier designed for riders can create inappropriate incentives for drivers because trips consume different amounts of time during surge and non-surge periods. The authors develop a dynamic stochastic model and derive an incentive-compatible driver payment mechanism, then compare it with real marketplace data. They find that an additive surge mechanism can better align driver incentives than a multiplicative mechanism. The strength is the combination of analytical derivation and real-data validation; the main limitation for the present topic is that sports-facility supply is not mobile labor. Nevertheless, the study contributes an important design principle: dynamic prices should reflect the opportunity cost of scarce time capacity.

\subsubsection{Surge pricing and driver behavior}
Miao, Wang, and Tang investigate the causal effect of surge pricing on driver behavior using a quasi-experiment \cite{r4}. The problem is the lack of clean causal evidence about whether surge changes labor supply and earnings. Their methodology combines a difference-in-differences design with causal-forest analysis after staggered implementation of surge pricing in two Chinese cities. They report higher weekly driver revenue but also countervailing cherry-picking and competition effects; daily revenue decreased because competition dominated at that margin, while drivers worked more days. The study advances knowledge by separating intensive and extensive labor responses. Its quasi-experimental design is a strength, while the mobile-worker setting limits direct transfer to fixed sports facilities. The relevance is that price changes can alter behavior in several directions, so a sports marketplace should measure both booking volume and changes in booking timing.

\subsubsection{Dynamic discounts for shared rides}
Jiao and Ramezani study dynamic discounting to incentivize shared rides in e-hailing markets \cite{r5}. The problem is how a platform can use personalized, time- and space-dependent discounts to increase shared rides without sacrificing profit. The methodology develops a non-equilibrium model in which passengers differ in price sensitivity, service quality sensitivity, and impatience. The authors optimize discounts and show that dynamic discounts can increase ridesharing and platform profit while reducing the fleet needed to serve demand. The model explicitly represents heterogeneous customers, which is a strength. A limitation is that ride-sharing involves matching and mobility constraints absent from fixed sports venues. For sports facilities, the analogous idea is targeted off-peak discounts that move demand from low-utilization slots toward periods with available capacity.

\subsubsection{Reinforcement learning for dynamic pricing}
A 2023 study in Transportation Research Part B develops scalable reinforcement-learning approaches for ride-hailing dynamic pricing \cite{r6}. The problem is the computational difficulty of solving large dynamic-pricing problems. The authors formulate pricing as a sequential decision problem, distinguish historical from current prices, establish an optimal stationarypolicy result, and improve the TD3 reinforcement-learning algorithm with three mechanisms. Experiments on small and large problems show improved performance of the proposed approaches. The strength is scalability and algorithmic evaluation; a limitation is the need for careful training and state representation, which can be challenging with sparse booking histories. The method is relevant because a sports marketplace may eventually need a policy that reacts continuously to time, capacity, weather, and demand rather than recomputing prices with a static rule.

\subsubsection{Dynamic pricing with cancellation behavior}
Xu, Saberi, and Liu study dynamic pricing and penalty strategies when customers can cancel orders over time \cite{r7}. The problem is that time-dependent cancellations affect platform utilization and the profitability of price and penalty policies. Their methodology models both unconfirmed and confirmed order cancellations and evaluates alternative pricing and penalty regimes. The results show that relaxing some pricing or penalty restrictions can increase profit and taxi utilization under the modeled conditions. The strength is the explicit treatment of cancellation timing, which is often ignored in simple demand models. The limitation is the transportation setting and the dependence on behavioral assumptions. The relevance to sports booking is direct: cancellation and no-show probabilities should be included when forecasting effective occupancy and setting prices for scarce slots.

\subsubsection{Dynamic equilibrium and pricing under uncertainty}
Liang, Jiang, and Lo examine a ride-sourcing market under travel-time uncertainty \cite{r8}. The problem is how uncertain travel times affect departure choices, mode choice, market equilibrium, and platform pricing. Their methodology formulates a dynamic equilibrium model incorporating simultaneous departure-time and mode choices and then examines pricing strategies. The study shows how uncertainty changes equilibrium outcomes and identifies pricing strategies for the platform. A strength is the explicit treatment of uncertainty and behavioral equilibrium; a limitation is the transportationspecific structure. For sports facilities, the transferable insight is that uncertainty should enter both demand forecasting and pricing decisions rather than being treated only as forecast error.

\subsection{Forecasting and sports pricing applications}

\subsubsection{Recreational foot-traffic forecasting}
Madden et al. forecast daily foot traffic at recreational trails using machine learning \cite{r9}. The problem is to provide managers with forecasts that can support adaptive management of visitor levels. The study uses data from three walking trails and incorporates weather factors such as rain and wind. The results indicate satisfactory predictive accuracy and show that weather variables help explain visitation. The study's main contribution is an empirical demonstration that outdoor recreational demand can be forecast from environmental and contextual variables. Its strength is the direct connection between forecasting and a recreation setting; limitations include the daily rather than slotlevel horizon and the absence of pricing. This article is highly relevant because sports facilities, especially outdoor pitches and courts, face similar weather-sensitive demand.

\subsubsection{Sports ticket price determinants}
Quansah, Buraimo, and Langa examine determinants of matchday ticket prices in the English Premier League \cite{r10}. The problem is identifying which observable sporting and market factors explain differences in ticket prices. The methodology uses regression models for the cheapest and most expensive tickets across five seasons. Results show that opponents and local derbies affect prices, while team performance has an asymmetric effect concentrated on more expensive tickets. The study contributes empirical evidence that sports prices reflect both event characteristics and market demand. Its strength is multiseason real-world sports data; its limitation is that tickets for discrete events differ from repeated bookable time slots. The relevance is that sports-specific demand signals, including opponent/event characteristics, can complement generic time-series features in a sports marketplace.

\subsubsection{Centralized pricing in online marketplaces}
Filippas, Jagabathula, and Sundararajan study a sharing-economy marketplace that moved from decentralized to centralized pricing \cite{r11}. The problem is how centralized platform pricing changes asset utilization, revenue, and provider behavior. The authors use experimental and quasi-experimental evidence from the marketplace. Centralized pricing increased provider-asset utilization and revenue, but providers without pricing control made non-price adjustments, including reducing availability and exiting; partial pricing control produced similar revenue gains with fewer reactions. The study's major contribution is empirical evidence that pricing authority itself changes supply behavior. The strength is real marketplace evidence; a limitation is the specific platform context. For sports facilities, it implies that owners may need bounded control over platformgenerated prices rather than fully opaque centralized pricing.

\subsubsection{Dynamic pricing with external information and inventory}
Li and Zheng study dynamic pricing when a seller observes new external information while holding finite inventory and does not know the true demand function \cite{r12}. The problem combines demand learning with inventory management. The authors formulate online and clairvoyant regimes and derive bounds on the revenue gap between them, then develop dualbased dynamic-pricing algorithms. The work contributes a rigorous framework for learning demand while making revenue decisions under capacity constraints. Its strengths are explicit uncertainty, inventory, and learning; a limitation is the stylized single-seller formulation. The relevance is strong because a sports slot is exactly a finite inventory unit whose value changes as the booking deadline approaches and as new weather or demand information arrives.

\subsubsection{Dynamic pricing with reviews}
Shin, Vaccari, and Zeevi analyze how online reviews change the demand function and therefore the optimal dynamic price \cite{r13}. The problem is that customers learn product quality from previous reviews, so willingness to pay changes endogenously over time. The methodology uses quality-based and value-based review models, a fluid approximation, and simulation using marketplace data. The study identifies time-scale regimes in which dynamic pricing differs from fixed pricing and shows that the review model affects the structure of optimal prices. Strengths include explicit social learning and a tractable dynamic model; the limitation is the assumption of a monopolist selling a product rather than a capacity-constrained sports platform. For sports facilities, the implication is that ratings and reviews can become an additional state variable in price optimization.

\subsubsection{Dynamic pricing with ratings}
Carnehl, Stenzel, and Schmidt study dynamic pricing when prices affect online ratings and ratings subsequently affect demand \cite{r14}. The problem is a feedback loop: high prices may reduce perceived value and worsen reviews, but may also select customers with stronger preferences who give better reviews. The methodology is an analytical model of dynamic pricing with endogenous ratings and customer selection. The results characterize optimal pricing and show that the platform's preferred rating sensitivity depends on how price affects reviews and how the platform weighs consumer and seller objectives. The strength is the explicit feedback mechanism; the limitation is the stylized product-market setting. For sports facilities, reviews can similarly interact with price, perceived quality, and future bookings, so a mature system could incorporate rating dynamics.

\subsection{Forecasting and deep learning reviews}

\subsubsection{Forecasting theory and practice}
Petropoulos et al. provide a broad review of forecasting theory and practice across methods, principles, evaluation, and applications \cite{r15}. The scope is intentionally wide rather than limited to one industry. Major themes include forecast preparation, model selection, evaluation, uncertainty, and real-world applications. The review emphasizes that different applications require different forecasting approaches and that forecast evaluation is part of the forecasting process. Its main future direction is continued development of methods and better integration of forecasting principles with applications. The paper helps shape this project by supporting a model-comparison strategy rather than assuming that one algorithm is universally superior.

\subsubsection{Deep learning for time-series forecasting}
Benidis et al. survey deep-learning approaches for time-series forecasting and organize the field around neural forecasting architectures, probabilistic forecasting, covariates, and large-scale industrial applications \cite{r16}. The review reports that deep learning has become prominent in forecasting and frequently performs strongly in forecasting competitions. Common approaches include recurrent architectures, temporal convolutional models, attention-based models, and probabilistic methods. The review highlights challenges involving data requirements, interpretability, uncertainty, and model selection. Its contribution to the present topic is methodological: a sports platform can compare classical statistical baselines with deep models instead of deploying a complex model without evidence that it improves short-horizon slot forecasts.

\subsubsection{Occupancy prediction by machine learning}
Zhang, Wu, and Calautit review 160 machine-learning studies on occupancy prediction in buildings \cite{r17}. The scope covers occupancy behavior, data collection, algorithms, and validation. The review identifies increasing use of neural-network methods and evaluates the workflow from data acquisition through prediction and validation. Gaps include the need for better data collection and reliable validation across contexts. Common methodologies include neural networks and other machine-learning models using temporal and environmental variables. Although building occupancy differs from sports bookings, the paper provides a useful conceptual analogy: both problems predict presence or utilization from historical patterns and contextual variables.

\subsubsection{Comprehensive review of occupancy prediction}
Li et al. systematically review 263 studies on building occupancy prediction \cite{r18}. The scope covers data-collection technologies, predictive algorithms, validation metrics, and their interaction. The review identifies Internet-of-Things data collection and hybrid machine-learning models as important directions and emphasizes that model selection should be linked to available data. It also identifies continuing needs for improved data collection and algorithm development. The dominant methodologies include machine learning, hybrid models, and sensor/IoT data. For sports facilities, the lesson is that forecasting quality depends not only on the algorithm but also on the richness and reliability of booking, weather, calendar, and venue-level data.

\subsubsection{Long sequence time series forecasting}
Chen et al. survey long-sequence time-series forecasting with deep learning and propose a taxonomy for long-horizon forecasting methods \cite{r19}. The review distinguishes different interpretations of longsequence forecasting, compares methods, and discusses evaluation resources and future research directions. Transformer and other deep-learning architectures are central themes, while the authors also emphasize standardized evaluation. The gap most relevant to this project is the difficulty of generalizing model performance across datasets and horizons. The paper helps define an experimental protocol: sports-slot forecasting should be evaluated with consistent metrics and time horizons and compared against strong baselines.

\subsubsection{Transformer-based long-term forecasting}
Su et al. provide a systematic review of transformer-based long-term time-series forecasting \cite{r20}. The scope covers the evolution of transformer architectures, their use in forecasting, experimental comparisons, and future directions. The review shows the rapid adoption of attention-based architectures but also points to efficiency, generalization, and data-related challenges. The common methodology is comparative evaluation of transformer variants on benchmark forecasting tasks. The main contribution to this project is identifying transformers as a candidate advanced model, while the limitations caution against assuming that their complexity is justified for small venue-level datasets.

\subsubsection{Deep learning-based time-series forecasting}
Song et al. review deep-learning time-series forecasting from 2014 to 2024 \cite{r21}. The review examines how models capture temporal and cross-variable correlations, how long-horizon computation can be reduced, and how different loss functions are used. It compares univariate and multivariate forecasting and discusses trends, seasonality, and computational efficiency. Future directions include better information extraction and more efficient forecasting architectures. For sports facilities, the review supports multivariate forecasting because weather, day-of-week, holidays, local events, price, and historical demand can jointly influence bookings.

\subsubsection{Demand forecasting under market volatility}
Muth, Lingenfelder, and Nufer systematically review machinelearning applications for demand prediction under macroeconomic volatility \cite{r22}. The review covers 64 studies and classifies publication trends, conceptual structures, and implementation strategies. It identifies the value of forward-looking data, economic indicators, coefficient-of-variation-based demand modelling, combined algorithms, and neural networks. A key gap is the need to connect machine learning with changing market conditions instead of treating demand as stationary. The review is useful for sports platforms because demand can also shift due to weather, holidays, local events, school calendars, and changing consumer behavior.

\subsubsection{Machine learning and deep learning for supply chain demand forecasting}
Mabrouki et al. critically review machine-learning and deep-learning models for demand forecasting in supply-chain management \cite{r23}. The scope compares forecasting algorithms and discusses their suitability for different supply-chain demand characteristics. Common approaches include classical machine learning, recurrent neural networks, and other deep-learning models. The review emphasizes that algorithm performance depends on data properties and application context. A future direction is better matching of algorithms to demand characteristics and operational constraints. The paper supports the proposed sports-facility study by reinforcing the need for baseline models, contextspecific evaluation, and careful feature engineering rather than selecting a model solely because it is newer.

\subsection{Platform and sharing economy reviews}

\subsubsection{Pricing strategies in the sharing economy}
Yang and Xia conduct a systematic literature review of 158 articles on pricing strategies in the sharing economy using Scopus and Web of Science \cite{r24}. The review classifies studies by business model and pricing party and identifies scenarios such as accommodation, ride sharing, ride hailing, bike sharing, vehicle sharing, crowdfunding, and crowdsourcing. Common methods include supply-demand balancing, optimization, auctions, game theory, equilibrium analysis, regression, and causal empirical methods. The authors identify gaps involving information asymmetry, competition, and the need for empirical approaches. This is one of the most directly relevant reviews because a sports facility aggregator is also a platform that coordinates underutilized capacity through pricing.

\subsection{Synthesis and research gap}
The studies identify several requirements for a sports-slot system. Forecasting must represent time, weather, calendar effects and venue differences \cite{r9,r15,r16,r17,r18,r19,r20,r21}. Pricing models must distinguish the effect of a price change from the demand conditions that led an owner to choose that price \cite{r2,r3,r4,r12}. Capacity expires at a known time, and cancellations affect the number of bookings that produce revenue \cite{r7,r12}. Facility owners also need control over price limits; the marketplace evidence connects pricing authority with changes in provider participation \cite{r11,r24}.

The supplied review identifies limited direct evidence for repeated bookable sports slots. Sports ticket studies concern discrete events, recreational visitation studies generally do not optimize price, and transport models often include mobile supply. The proposed study therefore combines booking forecasts, an explicit price-response model, cancellation prediction and a constrained slot-price optimizer. RQ1 evaluates the forecast component, RQ2 evaluates the pricing calculation, and RQ3 evaluates the operational outcomes. The simulation examines those components under a specified data-generating process; validation on real aggregator records remains a separate stage.
